\subsection{CONTRAGREDIENT REPRESENTATION}

Let \(E^{*}\) be the dual of E. Since W acts on E via \(\sigma\) , it also acts, by transport of structure, on \(E^{*}\) . The corresponding representation

\[
\sigma^ {*}: \mathrm{W} \to \mathbf {G L} (\mathrm{E} ^ {*})
\]

is called the contragredient representation of \(\sigma\) . We have

\[
\sigma^ {*} (w) = ^ {t} \sigma \left(w ^ {- 1}\right) \text {   for   all   } w \in W.
\]

If \(x^{*} \in \mathrm{E}^{*}\) and \(w \in \mathrm{W}\) , we denote by \(w(x^{*})\) the transform of \(x^{*}\) by \(\sigma^{*}(w)\) .

If \(s \in S\) , denote by \(A_s\) the set of \(x^* \in E^*\) such that \(x^*(e_s) > 0\) . Let \(C\) be the intersection of the \(A_s\) , \(s \in S\) . When \(S\) is finite, \(C\) is an open simplicial cone in \(E^*\) (§1, no. 6).

THEOREM 1 (Tits). If \(w \in W\) and \(C \cap w(C) \neq \emptyset\) , then \(w = 1\) .

We indicate immediately several consequences of this theorem:

COROLLARY 1. The group W acts simply-transitively on the set of \(w(\mathbf{C})\) for \(w \in \mathbf{W}\) .

This is clear.

COROLLARY 2. The representations \(\sigma\) and \(\sigma^{*}\) are injective.

Indeed, if \(\sigma^{*}(w) = 1\) , then \(w(\mathrm{C}) = \mathrm{C}\) , so \(w = 1\) by the theorem. The injectivity of \(\sigma\) follows from that of \(\sigma^{*}\) .

COROLLARY 3. If S is finite, \(\sigma(\mathrm{W})\) is a discrete subgroup of \(\mathbf{GL}(\mathrm{E})\) (provided with its canonical Lie group structure); similarly, \(\sigma^{*}(\mathrm{W})\) is a discrete subgroup of \(\mathbf{GL}(\mathrm{E}^{*})\) .

Let \(x^{*} \in C\) . The set U of \(g \in \mathbf{GL}(E^{*})\) such that \(g(x^{*}) \in C\) is a neighbourhood of the identity element in \(\mathbf{GL}(E^{*})\) ; by the theorem,

\[
\sigma^ {*} (\mathrm{W}) \cap \mathrm{U} = \{1 \};
\]

thus, \(\sigma^{*}(\mathbf{W})\) is a discrete subgroup of \(\mathbf{GL}(\mathbf{E}^{*})\) . By transport of structure, it follows that \(\sigma(\mathbf{W})\) is discrete in \(\mathbf{GL}(\mathbf{E})\) .

Proof of Theorem 1.

If \(w \in W\) , denote by \(l(w)\) the length of \(w\) with respect to S (Chap. IV, §1, no. 1).

We are going to prove the following assertions, where \(n\) denotes an integer \(\geqslant 0\) :

\((P_n)\) Let \(w \in W\) , with \(l(w) = n\) , and \(s \in S\) . Then:

either \(w(\mathbf{C})\subset \mathbf{A}_s\)

or \(w(\mathbf{C}) \subset s(\mathbf{A}_s)\) and \(l(sw) = l(w) - 1\) .

\((Q_n)\) Let \(w \in W\) , with \(l(w) = n\) , and \(s, s' \in S\) , \(s \neq s'\) . Let \(W_{s,s'}\) be the subgroup of \(W\) generated by \(s\) and \(s'\) . There exists \(u \in W_{s,s'}\) such that

\[
w (C) \subset u \left(A _ {s} \cap A _ {s ^ {\prime}}\right) \quad \text { and } \quad l (w) = l (u) + l \left(u ^ {- 1} w\right).
\]

These assertions are trivial for n = 0. We prove them by induction on n, according to the scheme

\[
\left(\left(P _ {n}\right) \text {   and   } \left(Q _ {n}\right)\right) \Longrightarrow \left(P _ {n + 1}\right) \quad \text { and } \quad \left(\left(P _ {n + 1}\right) \text {   and   } \left(Q _ {n}\right)\right) \Longrightarrow \left(Q _ {n + 1}\right).
\]

Proof that \(\left((\mathrm{P}_n)\text{ and } (\mathrm{Q}_n)\right) \Longrightarrow (\mathrm{P}_{n+1})\) .

Let \(w \in W\) , with \(l(w) = n + 1\) , and \(s \in S\) . We can write \(w\) in the form \(w = s'w'\) with \(s' \in S\) and \(l(w') = n\) . If \(s' = s\) , \((P_n)\) applied to \(w'\) shows that \(w'(C) \subset A_s\) , hence \(w(C) \subset s(A_s)\) and \(l(sw) = l(w') = l(w) - 1\) . If \(s' \neq s\) , \((Q_n)\) applied to \(w'\) shows that there exists \(u \in W_{s,s'}\) such that

\[
w ^ {\prime} (\mathrm{C}) \subset u \left(\mathrm{A} _ {s} \cap \mathrm{A} _ {s ^ {\prime}}\right) \text { and } l \left(w ^ {\prime}\right) = l (u) + l \left(u ^ {- 1} w ^ {\prime}\right).
\]

We have \(w(\mathbf{C}) = s'w'(\mathbf{C}) \subset s'u(\mathbf{A}_s \cap \mathbf{A}_{s'})\) .

Lemma 1. Let \(s, s' \in S\) , with \(s \neq s'\) , and let \(v \in W_{s,s'}\) . Then \(v(A_s \cap A_{s'})\) is contained in either \(A_s\) or in \(s(A_s)\) , and in the second case \(l(sv) = l(v) - 1\) .

The proof will be given in no. 5.

We apply the lemma to the element \(v = s'u\) . There are two possibilities: either

\[
  s ^ {\prime} u \left(\mathrm{A} _ {s} \cap \mathrm{A} _ {s ^ {\prime}}\right) \subset \mathrm{A} _ {s} \quad \text { and } a fortiori w (\mathrm{C}) \subset \mathrm{A} _ {s},
\]

or

\[
  s ^ {\prime} u \left(\mathrm{A} _ {s} \cap \mathrm{A} _ {s ^ {\prime}}\right) \subset s \left(\mathrm{A} _ {s}\right) \quad \text { and } \quad a fortiori w (\mathrm{C}) \subset s \left(\mathrm{A} _ {s}\right).
\]

Moreover, in the second case, \(l(ss'u) = l(s'u) - 1\) . Hence

\[
\begin{array}{c} l (s w) = l (s s ^ {\prime} w ^ {\prime}) = l (s s ^ {\prime} u. u ^ {- 1} w ^ {\prime}) \leqslant l (s s ^ {\prime} u) + l (u ^ {- 1} w ^ {\prime}) \\ = l (s ^ {\prime} u) + l (u ^ {- 1} w ^ {\prime}) - 1 \leqslant l (w) - 1, \end{array}
\]

and we know that this implies that \(l(sw) = l(w) - 1\) .

Proof that \(((\mathrm{P}_{n + 1})\) and \((\mathrm{Q}_n))\Longrightarrow (\mathrm{Q}_{n + 1}).\)

Let \(w \in W\) , with \(l(w) = n + 1\) , and \(s, s' \in S\) , \(s \neq s'\) . If \(w(C)\) is contained in \(A_s \cap A_{s'}\) , condition \((Q_{n+1})\) is satisfied with \(u = 1\) . For if not, suppose for example that \(w(C)\) is not contained in \(A_s\) . By \((P_{n+1})\) , \(w(C) \subset s(A_s)\) and \(l(sw) = n\) . By \((Q_n)\) , applied to \(sw\) , there exists \(v \in W_{s,s'}\) such that

\[
s w (C) \subset v \left(A _ {s} \cap A _ {s ^ {\prime}}\right) \quad \text { and } \quad l (s w) = l (v) + l \left(v ^ {- 1} s w\right).
\]

Then

\[
w (\mathrm{C}) \subset s v \left(\mathrm{A} _ {s} \cap \mathrm{A} _ {s ^ {\prime}}\right)
\]

and

\[
\begin{array}{c} l (w) = 1 + l (s w) = 1 + l (v) + l (v ^ {- 1} s w) \\ \geqslant l (s v) + l ((s v) ^ {- 1} w) \geqslant l (w), \end{array}
\]

so the inequalities above must be equalities. It follows that \((\mathbf{Q}_{n + 1})\) is satisfied with \(u = sv\) .

Proof of the theorem.

Let \(w \in W\) , with \(w \neq 1\) . We can write \(w\) in the form \(sw'\) with \(s \in S\) and \(l(w') = l(w) - 1\) . By \((\mathrm{P}_n)\) , applied to \(w'\) and \(n = l(w')\) , we have \(w'(C) \subset A_s\) , since the case \(w'(C) \subset s(A_s)\) is excluded because \(l(sw') = l(w) = l(w') + 1\) . So \(w(C) = sw'(C) \subset s(A_s)\) , and since \(A_s\) and \(s(A_s)\) are disjoint, we have \(C \cap w(C) = \varnothing\) . Q.E.D.

